\begin{frame}{Systems View on Deep Learning Models}
\framesubtitle{딥러닝 모델도 프로그램이다}

\begin{itemize}
  \item CNN, RNN, Transformer 등 딥러닝 모델의 구조는 다양한 종류가 있습니다.
  \item 딥러닝 모델은 다양한 구조를 가지지만, 중요한 특성을 공유합니다.
  \begin{itemize}
    \item \textbf{텐서 연산}으로 이루어져 있습니다.
    \item 전체적인 control flow가 대부분 한 방향만을 따라갑니다.
  \end{itemize}
  \pause
  \item 텐서: \only<2>{\(T: \underbrace{V^* \times \dots \times V^*}_{p\; \text{times}} \times \underbrace{V \times \dots \times V}_{q\; \text{times}} \to \mathbb{R}\) where \(V\) is a vector space ...}\only<3->{다차원 배열}
\end{itemize}
\only<4->{
\begin{figure}
  \centering
  \begin{subfigure}[c]{0.2\textwidth}
    {\resizebox{\linewidth}{!}{\definecolor{cA}{RGB}{66,133,244}
\definecolor{cB}{RGB}{234,67,53}
\definecolor{cC}{RGB}{52,168,83}
\definecolor{cGray}{RGB}{120,120,120}

\begin{tikzpicture}[
  font=\small,
  lbl/.style={font=\footnotesize\itshape},
  op/.style={font=\large},
]
  % ===== Geometry =====
  % cell = 0.3 units.
  % C: cols x 2.3..4.1 (6 cols), rows y 1.0..2.8 (6 rows)
  % Highlighted row: y 1.9..2.2 (4th row). Highlighted col: x 3.2..3.5 (4th col).
 
  % --- A: M x K, 2 cols, 6 rows. x 1.0..1.6 (2 cols), y 1.0..2.8 ---
  \draw[fill=cA!25,draw=cA] (1.0,1.0) rectangle (1.6,2.8);
  \draw[fill=cA!70,draw=cA] (1.0,1.9) rectangle (1.6,2.2);              % highlighted row
  \foreach \y in {1.3,1.6,1.9,2.2,2.5}{\draw[cA!60] (1.0,\y)--(1.6,\y);}
  \draw[cA!60] (1.3,1.0)--(1.3,2.8);                                   % 1 internal vertical
  \draw[cA] (1.0,1.0) rectangle (1.6,2.8);
  \node[lbl] at (1.3,0.65) {$A \in \mathbb{R}^{M \times K}$};
 
  % --- B: K x N, 2 rows, 6 cols. x 2.3..4.1, y 2.95..3.55 ---
  \draw[fill=cB!25,draw=cB] (2.3,2.95) rectangle (4.1,3.55);
  \draw[fill=cB!70,draw=cB] (3.2,2.95) rectangle (3.5,3.55);            % highlighted col
  \foreach \x in {2.6,2.9,3.2,3.5,3.8}{\draw[cB!60] (\x,2.95)--(\x,3.55);}
  \draw[cB!60] (2.3,3.25)--(4.1,3.25);                                 % 1 internal horizontal
  \draw[cB] (2.3,2.95) rectangle (4.1,3.55);
  \node[lbl,anchor=east] at (2.25,3.25) {$B \in \mathbb{R}^{K \times N}$};
 
  % --- C: M x N, 6 cols, 6 rows. x 2.3..4.1, y 1.0..2.8 ---
  \draw[fill=cC!25,draw=cC] (2.3,1.0) rectangle (4.1,2.8);
  \draw[fill=cC!80,draw=cC] (3.2,1.9) rectangle (3.5,2.2);             % output cell
  \foreach \x in {2.6,2.9,3.2,3.5,3.8}{\draw[cC!55] (\x,1.0)--(\x,2.8);}
  \foreach \y in {1.3,1.6,1.9,2.2,2.5}{\draw[cC!55] (2.3,\y)--(4.1,\y);}
  \draw[cC] (2.3,1.0) rectangle (4.1,2.8);
  \node[lbl] at (3.2,0.65) {$C \in \mathbb{R}^{M \times N}$};
 
  % arrows: A-row -> C-cell ; B-col -> C-cell
  \draw[-{Latex},cGray,thick] (1.6,2.05) -- (3.18,2.05);
  \draw[-{Latex},cGray,thick] (3.35,2.95) -- (3.35,2.22);
 
  % ===== Title =====
  \node[op] at (2.25, 4.05) {Matrix Multiplication};
 
  % ===== Normalize height for side-by-side placement =====
  \coordinate (Lx) at (current bounding box.west);
  \coordinate (Rx) at (current bounding box.east);
  \useasboundingbox (Lx |- 0,-0.2476) rectangle (Rx |- 0,4.7045);
\end{tikzpicture}}}
  \end{subfigure}%
  \begin{subfigure}[c]{0.3\textwidth}
    {\resizebox{\linewidth}{!}{\definecolor{cA}{RGB}{66,133,244}
\definecolor{cC}{RGB}{52,168,83}
\definecolor{cGray}{RGB}{120,120,120}

\begin{tikzpicture}[
  font=\small,
  lbl/.style={font=\footnotesize\itshape},
  op/.style={font=\large},
]
  % ---- grid drawing helper (3x3) ----
  \def\grid#1#2#3{% x y color
    \begin{scope}[shift={(#1,#2)}]
      \draw[fill=#3!22,draw=#3,thick] (0,0) rectangle (1.5,1.5);
      \foreach \k in {0.5,1.0}{
        \draw[#3!55] (\k,0)--(\k,1.5);
        \draw[#3!55] (0,\k)--(1.5,\k);}
    \end{scope}}
 
  % ===== Title & definition (centered on rendered image) =====
  \node[op]            at (2.889, 4.55) {Elementwise};
  \node[font=\footnotesize] at (2.889, 3.95) {$B = f(A_1, A_2, \dots, A_k)$};
 
  % ===== Operands: overlapping stacked squares =====
  % A_1 at top-left (back), stack moves down-right to A_k (front, on top)
  \grid{0.30}{2.00}{cA}
  \grid{0.55}{1.75}{cA}
  \grid{0.80}{1.50}{cA}
  % operand list below the stack
  \node[lbl] at (1.30,1.05) {$A_1, \dots, A_k \in \mathbb{R}^{N \times M}$};
 
  % ===== Arrow with f =====
  \draw[-{Latex},cGray,line width=1pt] (2.70,2.50) -- (3.90,2.50)
        node[midway,above,font=\itshape] {$f$};
 
  % ===== Result: single square =====
  \grid{4.30}{1.75}{cC}
  \node[lbl] at (5.05,1.30) {$B \in \mathbb{R}^{N \times M}$};
  %% ===== Normalize height for side-by-side placement =====
  \coordinate (Lx) at (current bounding box.west);
  \coordinate (Rx) at (current bounding box.east);
  \useasboundingbox (Lx |- 0,0.3226) rectangle (Rx |- 0,5.2748);
\end{tikzpicture}}}
  \end{subfigure}%
  \begin{subfigure}[c]{0.25\textwidth}
    {\resizebox{\linewidth}{!}{\definecolor{cA}{RGB}{66,133,244}
\definecolor{cC}{RGB}{52,168,83}
\definecolor{cGray}{RGB}{120,120,120}

\begin{tikzpicture}[
  font=\small,
  lbl/.style={font=\footnotesize\itshape},
  op/.style={font=\large},
]
  % Grids+arrow: A grid [0,1.6] (4x4), arrow, y grid [3.3,3.7] (1x4)
  % cell = 0.4. A: 4 cols x 4 rows. y: 1 col x 4 rows.
  \def\cx{1.9}
 
  % ===== Diagram =====
  % A: 4x4 grid, x 0..2.12, y 0..2.12, cell 0.53
  \begin{scope}[shift={(0.0,0.0)}]
    \draw[fill=cA!25,draw=cA] (0,0) rectangle (2.12,2.12);
    \draw[fill=cA!70,draw=cA] (0,1.06) rectangle (2.12,1.59);       % highlighted row
    \foreach \k in {0.53,1.06,1.59}{\draw[cA!60] (\k,0)--(\k,2.12);}
    \foreach \k in {0.53,1.06,1.59}{\draw[cA!60] (0,\k)--(2.12,\k);}
    \draw[cA] (0,0) rectangle (2.12,2.12);
  \end{scope}
  \node[lbl] at (1.06,-0.35) {$A \in \mathbb{R}^{M \times N}$};
 
  \draw[-{Latex},cGray,thick] (2.52,1.06) -- (3.42,1.06)
        node[midway,above,font=\footnotesize]{$\bigoplus_j$};
 
  % y: 4x1 grid, x 3.82..4.35, y 0..2.12
  \begin{scope}[shift={(3.82,0.0)}]
    \draw[fill=cC!25,draw=cC] (0,0) rectangle (0.53,2.12);
    \draw[fill=cC!80,draw=cC] (0,1.06) rectangle (0.53,1.59);       % highlighted entry
    \foreach \k in {0.53,1.06,1.59}{\draw[cC!60] (0,\k)--(0.53,\k);}
    \draw[cC] (0,0) rectangle (0.53,2.12);
  \end{scope}
  \node[lbl] at (4.085,-0.35) {$y \in \mathbb{R}^{M}$};
 
  % ===== Title & formula on grids+arrow center =====
  \node[op]            at (2.175, 3.05) {Reduction};
  \node[font=\footnotesize] at (2.175, 2.55) {$y_i=\textstyle\bigoplus_j A_{ij}$ (sum/max/\ldots)};
 
  % ===== Equal height + equal top/bottom padding (side-by-side) =====
  \coordinate (Lx) at (current bounding box.west);
  \coordinate (Rx) at (current bounding box.east);
  \useasboundingbox (Lx |- 0,-0.6588) rectangle (Rx |- 0,3.3563);
\end{tikzpicture}}}
  \end{subfigure}%
  \begin{subfigure}[c]{0.175\textwidth}
    {\resizebox{\linewidth}{!}{\definecolor{cA}{RGB}{66,133,244}
\definecolor{cC}{RGB}{52,168,83}
\definecolor{cGray}{RGB}{120,120,120}

\begin{tikzpicture}[
  font=\small,
  lbl/.style={font=\footnotesize\itshape},
  op/.style={font=\large},
]
  % Common vertical center axis
  \def\cx{2.7}
 
  % ===== Title & caption =====
  \node[op]            at (\cx, 2.90) {Reshape};
  \node[font=\footnotesize] at (\cx, 2.40) {same data, new shape};
 
  % source 2x6: 6 cols (width 2.52, center offset 1.05) -> shift_x = cx-1.05
  \begin{scope}[shift={(\cx-1.05,1.85)}]
    \foreach \r in {0,1}{
      \foreach \c in {0,...,5}{
        \pgfmathtruncatemacro{\val}{\r*6+\c+1}
        \node[draw=cA,fill=cA!20,minimum size=0.42cm,inner sep=0pt,font=\tiny]
          at (\c*0.42,-\r*0.42) {\val};}}
  \end{scope}
 
  % arrow (vertical, centered on cx, equal gaps to both grids)
  \draw[-{Latex},cGray,thick] (\cx,1.06) -- (\cx,0.70);
 
  % target 3x4: 4 cols (width 1.68, center offset 0.63) -> shift_x = cx-0.63
  \begin{scope}[shift={(\cx-0.63,0.35)}]
    \foreach \r in {0,1,2}{
      \foreach \c in {0,...,3}{
        \pgfmathtruncatemacro{\val}{\r*4+\c+1}
        \node[draw=cC,fill=cC!20,minimum size=0.42cm,inner sep=0pt,font=\tiny]
          at (\c*0.42,-\r*0.42) {\val};}}
  \end{scope}
  %% ===== Normalize height for side-by-side placement =====
  \coordinate (Lx) at (current bounding box.west);
  \coordinate (Rx) at (current bounding box.east);
  \useasboundingbox (Lx |- 0,-1.2273) rectangle (Rx |- 0,3.7249);
\end{tikzpicture}}}
  \end{subfigure}
\end{figure}}
\end{frame}